\subsection{Catálogo de interfaces}
\label{sec:catalogo-interfaces}

El catálogo identifica quince integraciones por número estable, dueño y superficie contractual. Los anexos 4-G y 4-H reúnen para cada una el modo, volumen esperado, ventana requerida de la contraparte y conducta ante lentitud o error (RT-05.21; Bases Técnicas Transversales, cap. 5, p. 13). Se distingue una ventana requerida por el diseño de un SLA efectivamente acordado con un tercero. El contrato ejecutable OpenAPI o AsyncAPI debe fijar operación o evento, esquema, versión, autenticación, errores y plazo de retirada. La publicación de una ruta genérica por sí sola no equivale a entregar ese contrato.

\subsubsection{Integraciones internas}

Las ocho interfaces internas intercambian pedidos, entregas, inventario, identidad y telemetría entre dispositivos, sitios y nube. El Anexo 4-G permite reconocer qué módulo acepta cada mensaje y cómo se conserva cuando falla el enlace.


INT-01 y INT-02 conservan el UUID de cada operación hasta obtener acuse durable; INT-03 y INT-04 ordenan eventos por sitio o agregado. INT-12 replica únicamente el WMS de Talca a un esquema de lectura separado de las tablas escritas por los consumidores de eventos. La extracción continua por fibra, LTE y satélite sostiene un desfase menor de 15 minutos; AL-DR-01 del Anexo 4-M verifica la pérdida del sitio y 4.3.2 justifica el límite residual.

\subsubsection{Integraciones externas}

Las siete interfaces con terceros se encapsulan para que un cambio de proveedor no altere doce módulos a la vez. El Anexo 4-H separa el tercero, el dueño interno del contrato y la conducta ante falla.

INT-06 traspasa la sugerencia de reposición aprobada de M2 al módulo de compras del ERP mediante la ACL (RF-10.03). El ERP conserva la emisión de las órdenes de compra (Subdocumento 3, D-05).


La integración tributaria INT-07 pasa por el ERP, único emisor de DTE; M5 no llama al SII. La falla de INT-09 impide afirmar que un cargo quedó aprobado. Para INT-08 se preservan el mensaje original, su equivalencia y la respuesta enviada a cada cadena.

INT-11 distribuye avisos por correo electrónico, SMS/WhatsApp y portal, separando los operacionales o transaccionales de las comunicaciones comerciales o promocionales. Los avisos necesarios para prestar el servicio, como confirmaciones de pedido, ETA e incidencias de entrega, no admiten baja comercial y no incorporan publicidad. Cada mensaje comercial identifica al remitente y su finalidad e incluye una vía expedita de baja: dirección válida y enlace en correo, enlace de suspensión en SMS/WhatsApp y control equivalente en portal. La baja registra destinatario, alcance, canal, fecha y origen en la preferencia de notificación; el trabajador Laravel consulta esa preferencia inmediatamente antes del envío y descarta también los mensajes comerciales pendientes o reintentados cuando existe una suspensión, que rige desde la solicitud. La aceptación comprueba la baja en cada canal y la continuidad de los avisos necesarios sin contenido promocional (Ley N.º 19.496, 1997, art. 28 B; RT-16.25; Bases Técnicas Transversales, cap. 16, p. 30).

Desde la consola Angular, el CLIENTE administra el cuerpo y la versión de las plantillas de avisos y comprobantes de pedido y entrega; Laravel combina la plantilla con los datos autorizados de la transacción y permite descargar el comprobante en HTML, formato abierto. Se reutilizan la plantilla de notificación, su cuerpo versionado y los tipos PEDIDO y ENTREGA, sin alterar la emisión tributaria del ERP; el editor admite texto y campos permitidos, sin ejecutar código de la plantilla (LafroX, 2026b, Anexo 5-A; RT-16.19; Bases Técnicas Transversales, cap. 16, p. 30).

\textbf{Volumen de mensajes por integración.} Los anexos 4-G y 4-H declaran las quince interfaces; el Anexo 4-I resume órdenes de magnitud y cálculos. No se usa la suma de eventos, consultas, muestras y reenvíos como número de transacciones únicas: una misma operación cruza varias interfaces. La prueba mide además tamaño de mensajes, objetos y WAL. La telemetría se ingiere y consolida por su flujo específico sin competir con las escrituras críticas de despacho (ADR-04).

